\todo{本章讲选了什么研究方法、为什么这样选（导师反馈第 1 条）。实验细节都在 Chapter~\ref{sec:implementation}。开头加一段 metatext。各节论点等中文思路。}

\subsection{Research Approach}\label{sec:methods:approach}

\todo{design science：构造一个 artifact（框架 + 原型），再用受控实验评估。为什么选这个方法，而不是纯文献研究或用户研究。}

\subsection{What Is Evaluated}\label{sec:eval:aim}

\todo{table: the framing table. Rows precision / recall; columns unconstrained action sources / this framework; cells measured, measured, ``guaranteed, not measured'', measured. Small and early: it defines the chapter.}

\todo{为什么测合法性（当前做不了的操作、用不了的控件），而不是一般的界面质量。}

\todo{Scope exclusions: interface aesthetics, completion time, user preference, and anything requiring human participants. 这里要写为什么不做用户研究。}

\subsection{Why a Synthetic Site}\label{sec:methods:synthetic}

\todo{A fully defined set of valid actions is the precondition for the dependent variable, which is why the site is synthetic. 和 Discussion 的 Limitations（Section~\ref{sec:disc:threats}）呼应。}

\subsection{An Independent Judge}\label{sec:methods:judge}

\todo{Why this is not circular: an independent behavioural judge that never reads the site's definition of validity. An action counts as available in $s$ iff some parameter assignment is not refused by the server. Same idea as the fail-closed validator in ASIL. 具体结果在 Section~\ref{sec:eval:method:truth}。}

\subsection{Overall Procedure}\label{sec:methods:procedure}

\todo{总体流程和 RQ 的对应：RQ1 文献比较（Chapter~\ref{sec:background}）；RQ2 设计并实现框架（Chapter~\ref{sec:implementation}）；RQ3 两个实验（Chapter~\ref{sec:results}）。可以用一张表。}
